% Prose rendering of PrimeTensor/Analysis/Growth.lean.
% Fragment: no \documentclass, no preamble, no document environment.

\section{Primitive response-scale growth control}
\label{Analysis:Growth:sec}

\leanfile{PrimeTensor/Analysis/Growth.lean}

\subsection*{What this file establishes}

\texttt{PrimeTensor/Analysis/Chain.lean} proved the chain rule under a
hypothesis, $\name{ChainAdmissibleAt}$, that quantifies over \emph{arbitrary
error germs}: it says that anything negligible against the inner response is
negligible against the perturbation.  That is a strong and rather indirect
thing to have to assume, and the module header names the defect precisely --- it
is a second-order transport assumption where a first-order geometric one ought
to do.

This file supplies the geometric condition.  $\name{ResponseScaleControlledAt}$
speaks only about the perturbation and its response, quantifying over no error
germs at all, and it implies admissibility.  The chain rule is then restated
with it in place of the transport hypothesis, and the two base cases --- the
identity and the constants --- are verified.

The interesting part is how a growth bound gets written down at all, since
$\name{Depth}$ has no subtraction.  The header states the technique: scale
comparison is expressed entirely through $\name{AtOrAfter}$.
Design note~\ref{Analysis:Growth:note:big-o} works out what that buys.

\subsection{The condition}

\begin{definition}[\lean{ResponseScaleControlledAt}]
\label{Analysis:Growth:def:rsc}
$f$ is \emph{response-scale controlled at} $x$ when, for every perturbation $h$
approaching the pivot and every target gain $t$, there are a response gain $r$
and an anchor such that, beyond the anchor and for every level $\ell$ the
perturbation achieves, the response achieves \emph{some} level $\rho$ with
\[
  \name{ScaleNear}(\rho,\; \name{response}(f,x,h)_n,\; 1)
  \quad\text{and}\quad
  \name{AtOrAfter}\bigl(\name{advance}(\ell,t),\; \name{advance}(\rho,r)\bigr).
\]
\end{definition}

\begin{leancode}
def ResponseScaleControlledAt
    (f : MulReal → MulReal) (x : MulReal) : Prop :=
  ∀ h : MulReal.Seq,
    ApproachesPivot h →
    ∀ targetGain : Depth,
      ∃ responseGain : Depth,
      ∃ anchor : Depth,
        ∀ n : Depth,
          Depth.AtOrAfter anchor n →
          ∀ level : Depth,
            MulReal.ScaleNear level (h n) 1 →
            ∃ responseLevel : Depth,
              MulReal.ScaleNear responseLevel
                (response f x h n) 1 ∧
              Depth.AtOrAfter
                (Depth.advance level targetGain)
                (Depth.advance responseLevel responseGain)
\end{leancode}

\begin{designnote}[A growth bound without subtraction]
\label{Analysis:Growth:note:big-o}
Read the concluding comparison numerically.  $\name{AtOrAfter}(a,b)$ is
$\size{a} \leq \size{b}$ and $\size{\name{advance}(\ell,g)} = \size{\ell} +
\size{g}$, so the condition is
\[
  \size{\ell} + \size{t} \;\leq\; \size{\rho} + \size{r},
  \qquad\text{that is}\qquad
  \size{\rho} \;\geq\; \size{\ell} + \size{t} - \size{r}.
\]
The response's level lags the perturbation's by at most the fixed amount
$\size{r} - \size{t}$.  In the logarithmic chart of
\texttt{PrimeTensor/Analysis/Derivative.lean}, where level $\ell$ is the
tolerance $2^{-(\size{\ell}-1)}$, that reads
\[
  \bigl\lvert \log \name{response}_n \bigr\rvert
    \;\leq\; 2^{\,\size{r}-\size{t}} \cdot
      \bigl\lvert \log h_n \bigr\rvert
    \quad\text{up to the dyadic rounding},
\]
which is the classical comparison $\Delta f = O(h)$ with
$2^{\size{r}-\size{t}}$ as the constant.

The inequality $\size{\rho} \geq \size{\ell} + \size{t} - \size{r}$ cannot be
stated directly: $\name{Depth}$ carries no subtraction, and none could be
added, since the depths start at one and a difference need not be positive.
The definition therefore never forms it.  It advances \emph{both} sides ---
the base level by $t$, the response level by $r$ --- and compares the two
results in the tail order of \texttt{PrimeTensor/Completion.lean}, which is
available.  Moving the deficit to the other side of the comparison turns a
subtraction into an addition, and addition of depths is what $\name{advance}$
already is.

This is worth setting against
\texttt{PrimeTensor/Analysis/Chain.lean}, which introduced admissibility as a
hypothesis on the ground that the required big-$O$ comparison had no expression
in the vocabulary.  That was true of the vocabulary as used there --- a single
level advanced by a gain, which can only be made finer --- but not of the
vocabulary as a whole.  Advancing both levels and comparing recovers exactly
the missing relation, which is what this file is for.
\end{designnote}

\begin{remark}[Where each quantifier sits]\label{Analysis:Growth:rem:quantifiers}
The response gain $r$ and the anchor are chosen after the perturbation and the
target gain but before the stage $n$ and the base level $\ell$, so a single $r$
must serve at every stage and every level --- the constant in the bound is
uniform.  The response level $\rho$, by contrast, is chosen after $\ell$, as it
must be, since how fine the response is depends on how fine the perturbation
is.  And the condition is stated per perturbation: a function may be
response-scale controlled with different constants along different $h$.
\end{remark}

\subsection{It implies admissibility}

\begin{theorem}[\lean{chainAdmissible\_of\_responseScaleControlled}]
\label{Analysis:Growth:thm:implies}
For every $f$ and $x$,
\[
  \name{ResponseScaleControlledAt}(f,x)
  \;\longrightarrow\;
  \name{ChainAdmissibleAt}(f,x).
\]
\end{theorem}

\begin{proof}
Let $h$ approach the pivot, let $\name{error}$ be negligible relative to the
response, and let a target gain $t$ be requested.  Take the response gain $r$
and the control anchor supplied by the hypothesis at $t$, take the anchor
supplied by negligibility \emph{at the gain $r$}, and join the two.

Beyond the join, suppose the base achieves level $\ell$ at stage $n$.  The
control produces a response level $\rho$ at which the response is near the
pivot, together with the comparison
$\name{AtOrAfter}(\name{advance}(\ell,t), \name{advance}(\rho,r))$.  Feeding
$\rho$ to the negligibility hypothesis --- legitimate, since the response is
near the pivot at $\rho$ --- puts $\name{error}_n$ near the pivot at
$\name{advance}(\rho,r)$.  The comparison says that level is at or after
$\name{advance}(\ell,t)$, so weakening once by $\name{scaleNear\_weaken}$ of
\texttt{PrimeTensor/Analysis/Control.lean} delivers the required
$\name{advance}(\ell,t)$.
\end{proof}

\begin{remark}[Monotonicity, spent a second time]\label{Analysis:Growth:rem:monotone}
The last step is the whole reason the comparison is stated in the tail order
rather than as an equation.  Having proved the error at some level, the proof
needs it at a possibly coarser one, and that is precisely the nesting of levels
established in \texttt{PrimeTensor/Analysis/Control.lean}.  The pattern is now
familiar: a bound is obtained at whatever level the machinery produces, and
monotonicity converts it to the level that was asked for.
\end{remark}

\subsection{The two base cases}

\begin{lemma}[\lean{responseScaleControlledAt\_identity} and
  \lean{responseScaleControlledAt\_constant}]
\label{Analysis:Growth:lem:examples}
The identity map, and every constant map, are response-scale controlled at
every point.
\end{lemma}

\begin{proof}
In both cases take the response gain equal to the target gain, the anchor
$\ctor{one}$, and the response level equal to the base level $\ell$; the
comparison is then between $\name{advance}(\ell,t)$ and itself, discharged by
$\ctor{here}$.

What differs is the first component.  For the identity the response at stage
$n$ is $\name{ratio}(x h_n, x) = h_n$, so it is near the pivot at $\ell$
exactly because the perturbation is --- the hypothesis is handed back
unchanged.  For a constant the response is $\name{ratio}(c,c) = 1$, near the
pivot at every level by reflexivity, and the hypothesis is discarded.
\end{proof}

The docstrings distinguish the two as \emph{exact} and \emph{maximal} control,
and the distinction is real even though the witnesses coincide: the identity
attains the bound with constant $1$ and could not do better, whereas the
constant map would satisfy the condition for any response gain whatever, its
response being the pivot on the nose.

\subsection{The chain rule, restated}

\begin{theorem}[\lean{hasMulDerivativeAt\_comp\_of\_responseScaleControlled}]
\label{Analysis:Growth:thm:chain}
Suppose $f$ has derivative $D_f$ at $x$ with $D_f$ scale controlled, $f$ is
response-scale controlled at $x$, and $g$ has derivative $D_g$ at $f(x)$ with
$D_g$ scale controlled.  Then $g \circ f$ has derivative
$\name{comp}(D_g,D_f)$ at $x$.
\end{theorem}

\begin{proof}
Convert the third hypothesis by Theorem~\ref{Analysis:Growth:thm:implies} and
apply the chain rule of \texttt{PrimeTensor/Analysis/Chain.lean}.
\end{proof}

\begin{remark}[What has changed and what has not]\label{Analysis:Growth:rem:what-changed}
All five hypotheses are now first-order conditions on $f$, $g$ and their
tangent maps: two are boundedness of a tangent map, two are differentiability,
and the fifth is a growth bound on a response.  None quantifies over error
germs.  That is the whole improvement, and it is a real one --- an assumption
about all possible errors has become an assumption about one specific germ.

The underlying proof is unchanged; $\name{ChainAdmissibleAt}$ remains the
notion the chain rule actually consumes, and this theorem is a wrapper.  No
converse is proved, so the two conditions are not known to be equivalent, and
nothing rules out an $f$ that is admissible without being response-scale
controlled.
\end{remark}

\begin{remark}[Still no nontrivial instance]\label{Analysis:Growth:rem:instances}
Response-scale control is verified for the identity and for constants and for
nothing else.  It is not shown closed under products, inverses, powers or
composition, and neither is $\name{ChainAdmissibleAt}$.  So
Theorem~\ref{Analysis:Growth:thm:chain} can still be discharged only where
\texttt{PrimeTensor/Analysis/Chain.lean} could --- with $f$ the identity or,
now, a constant --- and in both of those the composite's derivative was already
known from \texttt{PrimeTensor/Analysis/Examples.lean}.  The closure lemmas that
would make the rule usable are not in this file.
\end{remark}

\subsection*{Summary of the correspondence}

\begin{center}
\small
\begin{tabular}{@{}lll@{}}
\hline
Lean declaration & Kind & Here \\
\hline
\lean{ResponseScaleControlledAt}    & definition & Def.~\ref{Analysis:Growth:def:rsc} \\
\lean{chainAdmissible\_of\_responseScaleControlled} & theorem & Thm.~\ref{Analysis:Growth:thm:implies} \\
\lean{responseScaleControlledAt\_identity} & theorem & Lem.~\ref{Analysis:Growth:lem:examples} \\
\lean{responseScaleControlledAt\_constant} & theorem & Lem.~\ref{Analysis:Growth:lem:examples} \\
\lean{hasMulDerivativeAt\_comp\_of\_responseScaleControlled} & theorem & Thm.~\ref{Analysis:Growth:thm:chain} \\
\hline
\end{tabular}
\end{center}

\noindent
Every declaration of \texttt{PrimeTensor/Analysis/Growth.lean} appears above.
The file contains no \lean{sorry}, no axiom, and no unproved named proposition.
Design note~\ref{Analysis:Growth:note:big-o} and
Remarks~\ref{Analysis:Growth:rem:quantifiers},
\ref{Analysis:Growth:rem:monotone}, \ref{Analysis:Growth:rem:what-changed}
and~\ref{Analysis:Growth:rem:instances} are stated for the reader and are not
declarations of the file; the numerical and logarithmic readings of the
comparison are glosses and are not formalised.
